% ============================================================================
%  4. Montaje físico y normas de ensamble
% ============================================================================
\section{Montaje físico y normas de ensamble}
\label{sec:montaje}

Cada regla de montaje responde a un modo de falla concreto de los sistemas que
mezclan lógica y motores (tabla~\ref{tab:montaje}).

\begin{table}[H]
  \centering
  \caption{Normas de montaje físico y falla que previene cada una.}
  \label{tab:montaje}
  \small
  \begin{tabularx}{\textwidth}{@{}>{\raggedright\arraybackslash}p{2.4cm} L L@{}}
    \toprule
    \textbf{Regla} & \textbf{Implementación} & \textbf{Falla que evita} \\
    \midrule
    Masa común & GND de \SI{5}{\volt} y de Vmotor unidos en un solo punto cerca del DRV8833. &
      Sin referencia común el driver no interpreta los niveles lógicos del PIC; unirlas en varios puntos crea lazos de tierra por donde circula corriente de motor. \\
    Desacoplo del PIC & Cerámico de \SI{0.1}{\micro\farad} entre VDD y VSS (pines 11/12 y 32/31). &
      Picos de consumo del núcleo que causan reinicios por caída de VDD. \\
    Bus de potencia & Electrolítico de \SI{100}{\micro\farad} en paralelo con cerámico de \SI{0.1}{\micro\farad} entre VM y GND del driver. &
      El electrolítico entrega los picos de arranque e inversión; el cerámico filtra el ruido de alta frecuencia del PWM. \\
    Supresión en motores & Cerámico de \SI{0.1}{\micro\farad} entre terminales de cada motor. &
      Ruido de las escobillas y transitorios $dI/dt$ que se acoplan a la lógica. \\
    Regulación & LM7805 desde la batería para VDD = \SI{5}{\volt}. &
      La batería cae al arrancar los motores; el regulador mantiene estable la lógica. \\
    ICSP & MCLR, RB6 y RB7 libres hacia el conector del PICkit~3. &
      Permite reprogramar sin retirar el PIC del montaje. \\
    Reset & \SI{10}{\kilo\ohm} de MCLR a VDD y pulsador NA a GND. &
      MCLR flotante provoca reinicios aleatorios. \\
    Cableado & UTP calibre 24~AWG cortado a medida. &
      Cables largos actúan como antenas y elevan la inductancia de las rutas de corriente. \\
    \bottomrule
  \end{tabularx}
\end{table}

\subsection{Oscilador}

El cristal va lo más cerca posible de OSC1/OSC2, con dos capacitores a GND. La
hoja de datos del PIC16F877A recomienda 15 a \SI{33}{\pico\farad} para
cristales de 4 y \SI{20}{\mega\hertz}; el esquemático usa \SI{22}{\pico\farad}.
Los \SI{44}{\pico\farad} de la lista de materiales quedan fuera de ese rango y
pueden dificultar el arranque del oscilador.

\subsection{Nivel lógico del HC-05}

Su entrada RXD trabaja a \SI{3.3}{\volt}, por lo que la salida TX del PIC
(\SI{5}{\volt}) pasa por un divisor de \SI{1}{\kilo\ohm} y \SI{2}{\kilo\ohm}:
\begin{equation}
  V_{RXD} = \SI{5}{\volt}\cdot\frac{2}{1+2} \approx \SI{3.3}{\volt}
\end{equation}
En sentido contrario, el pin RX del PIC (entrada Schmitt trigger) requiere en
hoja de datos $0.8\,V_{DD} = \SI{4}{\volt}$ para leer un 1. Los
\SI{3.3}{\volt} del módulo suelen funcionar en la práctica; si aparecen errores
de recepción se agrega un convertidor de nivel.

\figurapendiente{Fotografía del montaje en protoboard}{fig:montaje}
